\chapter{Примитивы, на которых всё стоит}

Эта глава для читателя, который писал большие системы, но не строил распределённых баз данных. Каждый раздел объясняет один примитив и говорит, зачем он нужен именно нам. Если примитив вам знаком, раздел можно пропустить; в остальных главах на них есть ссылки.

\section{Реплицированный журнал и порядок}

Самая простая распределённая структура данных это журнал: последовательность записей, к которой можно только дописывать. Если несколько машин держат одинаковую копию журнала, они держат одинаковое состояние, потому что состояние есть результат применения записей по порядку. Все сложные вещи в распределённых базах сводятся к одному вопросу: как договориться о порядке записей.

Договориться можно двумя способами. Первый: назначить одну машину секвенсером, который присваивает записям номера, а остальные лишь проверяют и подписывают. Второй: провести консенсус, где порядок фиксируется голосованием без единого назначающего. Первый способ проще и живее, но секвенсер может остановиться или соврать. Второй сложнее, но не имеет единой точки.

\begin{likbez}[Зачем нам]
Журнал протокола есть ровно такой реплицированный журнал: записи о целях, задачах, назначениях, ревью и интеграциях. На уровнях L0 и L1 порядок задаёт секвенсер, а кор-ноды подписывают; на L2 порядок фиксируется консенсусом. Разделение порядка и подписи нужно, чтобы журнал продвигался, даже когда кор-ноды это ноутбуки, которые бывают выключены: секвенсер это всегда включённый сервис проекта, а подписи приходят, когда ноутбуки в сети.
\end{likbez}

\section{Отказы: сбой и византийское поведение}

Распределённые протоколы различают два класса отказов. Сбой: машина остановилась или отвалилась от сети, но не врёт. Византийское поведение: машина делает что угодно, включая ложь, противоречивые сообщения разным собеседникам и сговор с другими. Название идёт из задачи о византийских генералах, которые должны договориться об атаке, когда часть из них предатели.

Протоколы для сбоев (Raft, Paxos) требуют, чтобы большинство машин было живо. Протоколы для византийского поведения (PBFT, Tendermint, HotStuff) требуют, чтобы злонамеренных было меньше трети.

\begin{likbez}[Откуда берётся треть]
Пусть машин $n$, злонамеренных не больше $f$. Решение принимается, когда собрано $q$ подписей. Мы не можем ждать всех, потому что $f$ машин могут просто молчать; значит, $q \leq n - f$. Среди собранных $q$ подписей до $f$ могут быть от лжецов, которые в другом месте подписали противоположное. Чтобы два противоречивых решения не могли оба набрать $q$, их честные подписанты должны пересекаться: $2(q - f) > n - f$, то есть $q > (n+f)/2$. Соединяя с $q \leq n - f$, получаем $n > 3f$, откуда $n \geq 3f + 1$ и $q = 2f + 1$. При четырёх машинах терпим одну лживую, при семи две.
\end{likbez}

\begin{likbez}[Зачем нам]
Кор-ноды на L2 могут быть скомпрометированы: взломанная машина, отравленный инжектом агент, недобросовестный участник. Требование $n \geq 3f+1$ и кворум $2f+1$ гарантируют, что ни одна запись не будет совершена дважды с противоречивым содержимым и что $f$ лжецов не подделают историю. На L0 и L1 византийская устойчивость не заявляется, и кворум проще.
\end{likbez}

\section{Кворум и сертификат кворума}

Кворум это достаточное число подписей под одним и тем же содержимым. Набор подписей, собранный под записью, называют сертификатом кворума, QC. Полезное свойство: QC можно передавать и проверять офлайн. Тот, кто получил запись с QC, не должен ни с кем связываться, чтобы убедиться, что запись совершена: достаточно проверить подписи по известным ключам.

\begin{likbez}[Зачем нам]
Среда исполнения проверяет QC на записи назначения, прежде чем открыть санбокс, и на записи об интеграции, прежде чем продвинуть основную ветку. Ей не нужно доверять ни оркестратору, ни секвенсеру: подписи говорят сами за себя. Определение кворума и есть параметр, который меняется между уровнями L0, L1 и L2.
\end{likbez}

\begin{likbez}[Чего порог большинства не даёт]
Пусть подписантов трое, порог два, и один из них, $A$, злонамерен. $A$ подписывает две несовместимые записи на одно и то же место в журнале; честный $B$ подписывает первую, честный $C$ вторую. Обе записи получают сертификат большинства: $\{A, B\}$ и $\{A, C\}$. Порог защитил от одиночного решения, но не от раздвоения истории. Пересечение двух кворумов размера $t$ содержит не меньше $2t - n$ участников, и чтобы в нём гарантированно был честный при $f$ злонамеренных, нужно $2t - n > f$; для $n = 3$, $t = 2$, $f = 1$ это не выполнено. Поэтому L1 у нас не обещает единственности истории силой подписей: порядок там задаёт доверенный секвенсер, а подписи контролируют авторизацию. Единственность истории при лжецах даёт только L2 с полноценным византийским консенсусом.
\end{likbez}

\section{Лидер, аренда и fencing-токен}

Во многих системах есть роль, которую в каждый момент должна держать ровно одна машина: лидер, владелец, оркестратор. Оформляют её арендой: записью о том, что машина X держит роль до момента T, и продлевают heartbeat. Если heartbeat пропал, аренда истекает и роль переходит другому.

Аренда сама по себе не защищает от классической ошибки. Машина X держит аренду, зависает на минуту (сборка мусора, сеть, сон ноутбука), аренда истекает, роль переходит Y, а потом X просыпается и, ничего не зная, продолжает работать как лидер. Теперь лидеров два.

Лекарство называется fencing-токен: с каждой арендой выдаётся номер, который только растёт, и каждое действие лидера несёт этот номер. Тот, кто исполняет действия, помнит наибольший виденный номер и отвергает всё, что пришло с меньшим. Проснувшийся X действует с устаревшим номером, и его действия отбрасываются.

\begin{likbez}[Зачем нам]
Оркестратор есть арендуемая роль. Каждая его запись несёт эпоху и fencing-токен аренды; валидаторы отвергают записи с устаревшим токеном, но сохраняют их в ветке расхождений для аудита. Так двое оркестраторов не могут одновременно раздавать задачи, даже если один из них об этом не знает. Тот же приём защищает передачу сессии: принудительный захват отсекает предшественника токеном.
\end{likbez}

\section{Хеш-цепочка, дерево Меркла и доказательство включения}

Хеш-цепочка: каждая запись содержит хеш предыдущей. Изменить запись в середине нельзя, не изменив все последующие хеши, а последний хеш, голову, все знают. Дерево Меркла: над хешами записей строится двоичное дерево, где каждый узел есть хеш двух детей, и корень дерева фиксирует весь журнал одним числом. Полезное свойство дерева: чтобы доказать, что запись входит в журнал с данным корнем, достаточно показать логарифм от длины журнала хешей, а не весь журнал. Это доказательство включения.

Снапшот это подписанный набор из корня, длины журнала, головы и версии. Кто знает снапшот, тот может проверить любую запись по доказательству включения, не имея журнала целиком.

\begin{likbez}[Зачем нам]
Среда исполнения есть светлый клиент журнала: ей не нужно хранить журнал целиком, достаточно последнего снапшота с QC и доказательства включения для конкретной записи назначения или интеграции. Журналы прозрачности, на которых стоит инфраструктура сертификатов и подписи артефактов, устроены ровно так.
\end{likbez}

\section{Защита от отката}

Снапшоты подписаны, но подпись не мешает подсунуть светлому клиенту старый, тоже подписанный снапшот. Если клиент примет старый снапшот, он согласится с состоянием, в котором ещё не было, скажем, отзыва ключа или остановки. Защита: клиент хранит наивысшую принятую версию журнала как долговременное состояние безопасности, а не как кэш, и отвергает всё, что старше. Это состояние переживает перезапуски и не сбрасывается при ротации ключей.

\begin{likbez}[Зачем нам]
Без защиты от отката сертификат экстренной остановки можно было бы обойти, подсунув среде снапшот до остановки. Инвариант I8 требует высшей принятой версии именно поэтому.
\end{likbez}

\section{Атомарная замена и очередь интеграции}

Compare-and-swap: замени значение на новое, но только если текущее равно ожидаемому. Если кто-то успел изменить значение раньше, операция не выполняется, и надо повторить с новым ожиданием. Это единственный способ безопасно продвигать одну общую ссылку, за которую конкурируют несколько писателей.

В системах интеграции кода из этого выросло правило: никогда не сливать состояние, которое не было собрано и протестировано ровно в том виде, в каком оно попадёт в основную ветку. Очередь интеграции берёт кандидатов по одному, строит слияние с текущей головой, тестирует именно его и только потом продвигает голову.

\begin{likbez}[Зачем нам]
Одобрение и проверки привязаны к хешу. Если ревьюили коммит $R$, а в основную ветку попадает слияние $R$ с ушедшей вперёд головой $M$, то попадает другой хеш и, возможно, другое поведение, даже если сам патч не изменился. Поэтому среда строит интеграционного кандидата $I$, финальное одобрение и проверки относятся к $I$, а голова продвигается атомарно с $M$ на $I$; если голова уже не $M$, кандидат перестраивается, и одобрение с проверками повторяются. Это инвариант I3.
\end{likbez}

\section{Доступность данных}

Журнал хранит хеши тяжёлых объектов, а не сами объекты: тела задач, находки ревью, снимки сессий. Хеш гарантирует целостность объекта, но не его наличие. Кор-ноды могут совершить запись с хешем объекта, который потом никто не сможет получить, потому что единственная копия лежала на выключенном ноутбуке.

\begin{likbez}[Зачем нам]
Перед включением хеша в журнал требуется сертификат доступности: несколько независимых реплик подтвердили, что хранят объект. Для чекпоинтов сессии оркестратора это критично: без чекпоинта преемник стартует холодно.
\end{likbez}

\section{Случайность, которую никто не контролирует}

Если ревьювера выбирать случайно, кто-то должен сгенерировать случайность. Если её генерирует оркестратор, он может подобрать удобную. Если она известна заранее, оркестратор может подобрать момент или состав кандидатов, чтобы случайность привела к нужному ревьюверу; это называется grinding.

Решения: случайность собирается из вкладов нескольких сторон так, что ни одна не контролирует итог (пороговый маяк, commit-reveal), либо вычисляется проверяемой случайной функцией из секрета одной стороны так, что результат нельзя подобрать, а проверить может любой. И в любом случае порядок такой: сначала фиксируется множество кандидатов и объект, для которого выбирают, и только потом получается случайность.

\begin{likbez}[Зачем нам]
Множество кандидатов в ревьюверы замораживается записью до появления хеша ревизии, случайность получается после, и ревьювер вычисляется из идентификатора задачи, хеша ревизии и случайности. Любой валидатор пересчитывает выбор. Это половина инварианта I5.
\end{likbez}

\section{Сибил-атака}

В открытой системе один участник может завести много личностей и получить больше голосов, больше назначений, больше репутации. Единственная универсальная защита это стоимость личности. Доказано, что не существует нетривиальной репутационной системы, которая одновременно устойчива к множественным личностям и вознаграждает пропорционально вкладу; можно только сделать атаку дорогой.

\begin{likbez}[Зачем нам]
Нода это ёмкость, и ёмкость стоит денег: подписка, ключ API, железо для локальной модели. Это естественный барьер, но не абсолютный, поэтому он дополняется репутацией и адаптивной глубиной проверки. Публичный рой не обещает пропорциональной справедливости, он обещает дорогую атаку.
\end{likbez}

\section{Что мы берём у готовых протоколов}

Raft даёт нам модель снимка и лога для передачи состояния оркестратора: преемник загружает последний чекпоинт и докатывает записи журнала после него. PBFT, Tendermint и HotStuff дают модель византийского консенсуса для L2, включая смену лидера консенсуса при его молчании или лжи. Журналы прозрачности дают форму снапшотов и доказательств включения. Очереди интеграции дают правило никогда не сливать непроверенное. Ничего из этого мы не изобретаем; наш вклад в композиции и в том, что предметом журнала являются полномочия агентов.
